\chapter{Memory Management}
\label{ch:memory}

Every value your program works with lives somewhere in the computer's memory.
Most of the time Salam handles that for you and you need not think about it. But
the growable containers vectors, hash maps and a few advanced techniques do
involve memory you are responsible for. This chapter explains Salam's memory model
plainly, so you know which memory takes care of itself and which you must tend.

\section{Two kinds of memory}

It helps to picture memory in two parts. The \textbf{stack} holds your local
variables the \code{int}s, the \code{bool}s, the fixed-size arrays and structs.
This memory is managed completely automatically: when a variable comes into scope
it is created, and when its scope ends it is reclaimed. You have been relying on
this since Chapter~2 without a second thought, and you can continue to.

The \textbf{heap} holds things whose size is not known until the program runs, or
that must outlive the function that created them the growable buffer behind a
vector, the entries of a hash map. Heap memory must be explicitly released when no
longer needed, which is why those containers have a \code{free} method.

\section{Values that manage themselves}

The good news first: the great majority of what you write needs no memory
management at all.

\begin{itemize}[leftmargin=1.4em]
  \item \textbf{Scalars} \code{int}, \code{f64}, \code{bool}, \code{char}, an
        enum live on the stack and vanish automatically.
  \item \textbf{Fixed-size arrays and structs} live inline, on the stack, and are
        likewise reclaimed automatically when they go out of scope. Recall from
        Chapter~9 that assigning or passing a struct \emph{copies} it there is
        no shared hidden allocation to worry about.
  \item \textbf{Strings} are special: although a \code{str} points at heap bytes,
        Salam manages those bytes for you (see the box below). You never
        \code{free} a string.
\end{itemize}

\begin{deep}[In Depth: the string arena]
Joining strings with \code{+}, or calling functions like \code{str.Upper}, must
allocate room for the new text on the heap yet you never free a string. How?
Salam keeps a per-program \emph{arena} for string data: all the bytes that string
operations produce are allocated from it, and the whole arena is released at once
when the program ends. The trade-off is deliberate: you get the convenience of
never managing string memory, at the cost of strings living until the program
exits. For the short-lived, string-heavy programs that are most common, this is
exactly the right bargain.
\end{deep}

\section{Containers you must free}

The exception to ``it manages itself'' is the growable containers. A
\code{Vector<T>} or a \code{HashMap<K, V>} owns a heap buffer that it resizes as
it grows, and that buffer must be returned to the system when you are finished. So
every growable container has a \code{free} method, and you should call it:

\begin{salam}
func main:
    mut v : Vector<int> = Vector {}
    v.push(1)
    v.push(2)
    v.push(3)
    println "sum of three pushes:", v.get(0)[0] + v.get(1)[0] + v.get(2)[0]
    v.free()              // release the buffer
end
\end{salam}

\begin{progout}
sum of three pushes: 6
\end{progout}

This applies to \code{Vector}, \code{HashMap}, \code{Stack}, \code{Queue}, and the
other growable collections. The fixed-size array \code{int[5]}, by contrast, needs
no freeing it lives on the stack.

\subsection{Pairing creation with \texttt{free} using \texttt{defer}}

The cleanest way never to forget a \code{free} is to schedule it the moment you
create the container, with the \code{defer} you learned in Chapter~14. Written
together on adjacent lines, the release is guaranteed and impossible to overlook:

\begin{salam}
func sum_first_n(n : int) : int:
    mut v : Vector<int> = Vector {}
    defer v.free()                 // guaranteed cleanup, however we return
    for mut i = 1, i <= n, i = i + 1:
        v.push(i)
    end
    mut total : int = 0
    mut j : int = 0
    while j < v.len():
        total += v.get(j)[0]
        j += 1
    end
    ret total
end

func main:
    println sum_first_n(5)
end
\end{salam}

\begin{progout}
15
\end{progout}

Because \code{defer v.free()} runs whenever \code{sum\_first\_n} returns through
the normal path or any early \code{ret} the vector's memory is always released.
This ``create, defer free, use'' pattern is the single most important habit for
leak-free Salam code.

\begin{warn}
Two mistakes to avoid. \textbf{Forgetting to free} a container leaks its
memory harmless in a tiny script that exits immediately, but a real problem in a
long-running program that creates containers in a loop. \textbf{Using a container
after freeing it} is worse: the buffer is gone, and reading it is undefined
behaviour. Free a container exactly once, as the last thing you do with it ---
which is precisely what \code{defer} at the point of creation guarantees.
\end{warn}

\section{No hidden ownership rules}

If you have heard of languages with elaborate ``ownership'' and ``borrowing''
systems, you can relax: Salam has none of that ceremony. There are no lifetime
annotations to write, no borrow checker to satisfy, no move semantics to learn.
The model is the simple one above stack values manage themselves, growable
containers get a \code{free} and the discipline is the single habit of pairing
each container with a \code{defer free}.

This simplicity is a deliberate design choice. It keeps the language approachable
and the code readable, at the cost of asking you to remember \code{free} for the
handful of containers that need it. In practice that is a small, learnable
responsibility, and \code{defer} reduces it to a reflex.

\section{Manual memory with \texttt{mem}}

For most programs the containers are all the heap memory you will ever touch. But
when you need raw control implementing your own data structure, or working at
the boundary with C (Chapter~19) the \code{mem} module gives you direct
allocation, the same primitives the standard library's containers are built on:

\begin{salam}
import "mem"

func main:
    p : void* = mem.Allocate(64 as u64)     // ask for 64 bytes
    if p == null:
        println "allocation failed"
        ret
    end
    println "got 64 bytes"
    mem.Free(p)                            // hand them back
    println "released"
end
\end{salam}

\begin{progout}
got 64 bytes
released
\end{progout}

\code{mem.Allocate(n)} requests \code{n} bytes and returns a \code{void*} (a raw,
untyped pointer), or \code{null} if the request fails which you should check.
\code{mem.Free(p)} returns the memory. The module also offers
\code{mem.Reallocate} to resize a block, \code{mem.Set} to fill bytes, and
\code{mem.Copy} to copy them.

\begin{warn}
Manual memory is powerful and unforgiving: a missing \code{Free} leaks, a double
\code{Free} corrupts, and using freed memory crashes. You will rarely need
\code{mem} directly the built-in containers cover almost every case. Reach for
it only when you are building something the standard library does not provide, and
when you do, guard every allocation with a \code{defer mem.Free} just as you would
a container.
\end{warn}

\begin{deep}[In Depth: tracking down leaks]
Because Salam compiles to C, the usual C tooling works on Salam programs. The
compiler can build with AddressSanitizer (the \code{--asan} flag, or the
\code{salam memcheck} command) to catch use-after-free and out-of-bounds errors
at run time, and the \code{mem} module exposes counters allocations made, frees
done, bytes still live that let a program audit its own memory use. If you ever
suspect a leak, these turn ``something is wrong'' into a precise location.
\end{deep}

\section{Chapter summary}

\begin{itemize}[leftmargin=1.4em]
  \item Scalars, fixed arrays, and structs live on the stack and are reclaimed
        automatically; you never manage them.
  \item Strings are heap-backed but managed by a per-program arena never
        \code{free} a string.
  \item Growable containers (\code{Vector}, \code{HashMap}, \code{Stack},
        \code{Queue}) own heap buffers and must be \code{free}d.
  \item Pair every container creation with \code{defer x.free()} to guarantee
        cleanup on every exit path.
  \item Salam has no ownership/borrowing system the model is deliberately
        simple.
  \item \code{mem} (\code{Allocate}, \code{Free}, \code{Reallocate}, \code{Set},
        \code{Copy}) provides raw manual memory for advanced use; guard it
        carefully.
\end{itemize}

\section{Exercises}

\exercise Write a function that builds a \code{Vector<str>} of three greetings,
prints them, and frees the vector using \code{defer}.

\exercise Take a function that creates a \code{HashMap}, returns early on some
condition, and uses \code{defer} to ensure the map is freed on both paths.

\exercise Explain in a comment why a fixed array \code{int[10]} needs no
\code{free} but a \code{Vector<int>} does.

\exercise Use \code{mem.Allocate} and \code{mem.Free} to grab and release 128
bytes, checking for a \code{null} result, and guard the free with \code{defer}.

\exercise Describe, in your own words, what the string arena does and one
situation where its ``strings live until the program ends'' behaviour could matter.
